\section{Introduction}
\label{sec:intro}

\emph{Security ceremonies} extend security protocols with the human users who take part in them~\cite{ellison2007ceremony}. A formal analysis of a ceremony has to account for what humans actually do, including their errors~\cite{bella2012layered,basin2016modeling,bella2022modelling,fila2024nothing}. The \XMen{} approach~\cite{sempreboni2020x,sempreboni2023mutation} models human errors as \emph{mutations} of a ceremony specification: a mutated specification describes a deviating human, its effect is matched and propagated to the other roles, and an unmodified verifier, Tamarin~\cite{meier2013tamarin,basin2025modeling}, analyses it as it would analyse the original. Sandu et al.~\SCP{} added \forget{}, a mutation for memory lapses. Forgotten information is not deleted from the human's knowledge, which stays monotone; it is recorded in a separate \Forget{} set, and for each send of the human the mutation decides whether the intended message can still be derived without forgotten hypotheses. Depending on the answer, the send is kept, \emph{replaced} by a substitute the human can still build, or \emph{skipped}; a lapse that affects an internal check \emph{neglects} the check. This finds attacks that arise from lapses rather than from cryptography, such as a user who has forgotten one bank's password and sends another bank's password instead~\SCP.

The decision of~\SCP{} is a syntactic test: a derivation is unusable if one of its hypotheses is blocked by the \Forget{} set. Three phenomena that matter for ceremonies fall outside it. First, people lose skills as well as values: a user may hold every value a step needs and still fail to carry it out, because what was lost is a way of proceeding, not a fact~\cite{ryle1949concept,cohen1980preserved}. Second, forgetting is often a failure of access rather than of storage: a cue, such as a password hint, can make a forgotten value accessible at one step without the value being learnt again~\cite{tulving1974cue,tulving1973encoding}. Third, the test produces a verdict but no reason, so an analyst who sees a \replace{} cannot ask which lapse caused it or whether a cue would have prevented it. None of these is a further element of the \Forget{} set. They are relations between pieces of information that pull in different directions: by default a known value can be used; a lapse is an exception to this default; a cue is an exception to the exception.

\paragraph{This paper.} We develop a \emph{defeasible reasoning framework} for \forget{}. For each human transition, the human's cognitive state is compiled into a theory of the ambiguity-blocking, team-defeat Defeasible Logic DL($\partial$)~\cite{nute1994defeasible,antoniou2001representation}. Possession ($\Known$) and usability ($\Usable$) are separate atoms; each normal derivation of a candidate message is a named conclusion ($\Via$); a lapse of a value, of an inference rule, or of a whole derivation is a defeater against that conclusion; and a cue is a higher-priority rule that overrules the defeater for one hypothesis while the lapse stays recorded. The mutation is read off the proof tags of the $\Usable$ atoms, and the proof is the explanation.

The contribution is not that \forget{} can be written in Defeasible Logic, but what the formulation lets us establish, check and extend. We characterise exactly what the compiled theories conclude and show that this rests on identifiable proof-theoretic choices: recovery from several lapses holds under team defeat and fails without it, while ambiguity blocking and propagation agree on every compiled theory. We prove that the framework conservatively extends~\SCP{} and strictly separates states the \Forget{} set cannot distinguish, that decisions stay sound when derivation search is cut off, and that usability is decidable in polynomial time unless an exact derivation is forgotten. The analysis also shows where the framework stops: in its theories a lapse can only defeat, and no rule concludes that a derivation is unusable. Adversarial cues, which interfere with a correct value, cross this line: usability stops being monotone in the cues present, and the choice of defeasible semantics decides whether a human facing conflicting cues acts on an alternative or withholds.

\paragraph{Contributions.}
\begin{enumerate}
\item A defeasible reasoning framework for \forget{} (\cref{sec:framework}): a cognitive state with procedural lapses and cues, its compilation into DL($\partial$), and a decision procedure with explanations and explicit coverage outcomes.
\item A formal analysis (\cref{sec:properties}): an exact characterisation of the compiled theories, its dependence on team defeat and independence from the ambiguity regime, conservative extension of~\SCP{} and two strict separations, soundness and partial completeness under bounded search, and a polynomial characterisation of usability.
\item An implementation in the \XMen~2.0 tool~\SCP{} and its evaluation (\cref{sec:implementation,sec:evaluation}): differential testing against an independent encoding of DL($\partial$), which exposed a team-defeat defect in that encoding; a comparison with~\SCP{} on 5{,}000 random transitions; scalability; and Tamarin analyses in which the new decisions change verification outcomes.
\item A preliminary, not implemented, model of adversarial cues (\cref{sec:discussion}).
\end{enumerate}
The \forget{} mutation, the \Forget{} set, its blocking cases and the published algorithm are due to~\SCP{} (\cref{sec:background}); the rest of the paper is new.
